% !PRINTED 329
\weberpage{329}

%<<CONT
wenn das obere Zeichen gilt, falls von den beiden Stammdiskriminanten $\Delta'$, $\Delta''$ wenigstens eine positiv ist, das untere, wenn beide negativ sind.

Setzen wir
\[
k = k'\Delta'' + k''\Delta',
\]
und lassen $k'$, $k''$ Restsysteme nach $\Delta'$, $\Delta''$ durchlaufen, so durchläuft $k$ ein Restsystem nach $\Delta = \Delta'\Delta''$.

Die Multiplikation der beiden Summen
\weq{4}{\overset{k'}{\Sigma}(\Delta',k')\,e^{-\frac{2s\pi ik'}{\Delta'}}, \quad \overset{k''}{\Sigma}(\Delta'',k'')\,e^{-\frac{2s\pi ik''}{\Delta''}}}
ergibt
\weq{5}{\overset{k'k''}{\Sigma}(\Delta',k')(\Delta'',k'')\,e^{-\frac{2s\pi ik}{\Delta}}.}

Es ist aber nach \S\,85, (16)
\[
(\Delta',k) = (\Delta',\Delta'')(\Delta',k'),
\]
\[
(\Delta'',k) = (\Delta'',\Delta')(\Delta'',k''),
\]
und folglich nach \S\,85, (13)
\[
(\Delta,k) = \pm(\Delta',k')(\Delta'',k''),
\]
und demnach ergibt sich aus (3), (4) und (5) die allgemeine Gültigkeit der Formel (2) für jede Stammdiskriminante $\Delta$.

\begin{center}\rule{0.2\linewidth}{0.4pt}\end{center}
